\chapter{PROFIL INDUSTRI}

\section{Profil Bolabot \textit{Techno Robotics Institute}}
Bolabot \textit{Techno Robotics Institute} merupakan institusi yang mengutamakan riset dan pendidikan di bidang teknologi robotika. Lembaga ini secara aktif melakukan penelitian yang mencakup otomasi perangkat, pengendalian alat, pemrosesan sinyal, penerapan kecerdasan buatan, sistem penglihatan untuk robot, serta instrumentasi dalam pengukuran alat-alat fisika \cite{15}.

Tidak hanya terbatas pada riset, institusi ini juga menyelenggarakan program pendidikan robotika bagi berbagai jenjang, mulai dari tingkat sekolah dasar hingga perguruan tinggi, serta menyediakan pelatihan khusus bagi para guru. Selain itu, Bolabot \textit{Techno Robotics Institute} mengembangkan beragam perangkat yang mendukung kegiatan praktikum pada mata pelajaran Elektronika Dasar, Fisika Dasar, dan Fisika Komputasi \cite{15}.

Di samping fungsinya sebagai perusahaan yang berfokus pada riset dan pendidikan, lembaga ini juga berhasil membentuk komunitas yang mengakomodasi antusiasme masyarakat terhadap robotika. Komunitas tersebut pertama kali didirikan pada tahun 2014 dengan nama \textit{“Physics of Instrumentation and Computation Organization”} (PICO), dan kemudian mengalami perubahan nama pada pertengahan tahun 2015 menjadi “Komunitas Pecinta Robot” (KOTARO) \cite{15}.

\section{Sejarah Bolabot \textit{Techno Robotic Institute}}
BOLABOT bermula dari komunitas riset di bidang komputasi dan instrumentasi yang dibentuk di lingkungan Fisika UIN Sunan Gunung Djati Bandung sejak tahun 2010. Inisiatif awal ini diprakarsai oleh beberapa tokoh kunci, yaitu Mada Sanjaya W.S., Dian Syah Maulana, Halimatussadiyah, Okyza Maherdi P, Irfan Syafar Farouk, Sandy Yusuf, Utin Sutinah, serta Aah Kholifah. Seiring berjalannya waktu, Irawan Febrianto sebagai alumni Instrumentasi UIN SGD angkatan 2008 dan 2009—bergabung untuk semakin memperkuat fondasi komunitas ini \cite{15}.

BOLABOT secara resmi mulai beroperasi pada tanggal 11 Maret 2012 di bawah payung hukum CV. Sanjaya Star Group, yang menandai dimulainya perjalanan strategis dalam bidang riset dan inovasi robotika. Kegiatan awal yang dilaksanakan meliputi penyelenggaraan program pembelajaran robotika secara reguler dan privat, pengembangan projek alat praktikum dalam mata pelajaran elektronika dasar, serta penyelenggaraan seminar dan workshop yang mendalami aspek teknologi robotika. Selanjutnya, upaya perluasan dampak pendidikan diwujudkan melalui pengenalan ekstrakurikuler robotika di berbagai sekolah. Pada bulan Desember 2012, BOLABOT bersama dengan Robotika UIN meluncurkan event Robocom Pertama, yang mencakup kompetisi robotik. Acara tersebut terbagi dalam dua segmen, yakni lomba robot \textit{line follower} analog bagi siswa binaan BOLABOT dan kompetisi \textit{robot soccer} digital yang melibatkan mahasiswa pada mata kuliah elektronika dasar Fisika Sains UIN Sunan Gunung Djati Bandung \cite{15}.

Lebih lanjut, pada tanggal 10 November 2014, terbentuklah komunitas khusus mahasiswa Instrumentasi Fisika UIN SGD dengan nama PICO \textit{(Physics Instrumentation and Computation Organization)}. Komunitas ini kemudian mengalami transformasi pada bulan Juli 2015 dan berganti nama menjadi KoTaRo (Komunitas Pecinta Robot), sebagai salah satu divisi sosial BOLABOT yang bertujuan untuk menyebarluaskan pengetahuan robotika secara lebih luas di Indonesia. Sejalan dengan dinamika strategis organisasi, pada tahun 2018 terjadi perubahan struktural di mana CV. Sanjaya Star Group secara resmi diubah menjadi CV. BOLABOT. Pada tahun yang sama, Divisi Publikasi Riset mendirikan penerbitan buku ilmiah yang dinamakan Penerbit BOLABOT. Inisiatif ini memungkinkan lembaga untuk menerbitkan karya ilmiah hasil riset internal maupun karya dari penulis lain, yang kemudian dapat dipasarkan secara komersial, memperkuat posisi BOLABOT sebagai pionir dalam riset dan inovasi di bidang robotika \cite{15}.

\section{Visi dan Misi Bolabot \textit{Techno Robotics Institute}}

\subsection{Visi Bolabot \textit{Techno Robotics Institute}}
Visi dari Bolabot \textit{Techno Robotics Institute} adalah menjadi ikon perusahaan berbasis edukasi, riset, dan teknologi robotika di Indonesia \cite{15}.

\subsection{Misi Bolabot \textit{Techno Robotics Institute}}
Misi dari Bolabot \textit{Techno Robotics Institute} adalah sebagai berikut \cite{6}:
\begin{enumerate}
    \item Mengembangkan pendidikan robotik yang menarik, kreatif, dan menyenangkan.
    \item Menyelenggarakan pendidikan nonformal sebagai penunjang kecakapan dan pembelajaran.
    \item Melakukan berbagai penelitian untuk menciptakan inovasi teknologi.
    \item Menciptakan berbagai produk untuk memenuhi kebutuhan pendidikan dan industri.
    \item Menciptakan usaha-usaha baru di bidang teknologi.
\end{enumerate}


\section{Roadmap Bolabot \textit{Techno Robotics Institute}}

Terdapat tiga bidang utama dan satu bidang sosial yang menjadi target kerja Bolabot, yaitu sebagai berikut \cite{15}:

\begin{enumerate}
    \item \textbf{Bidang Edukasi Robotika}: Menyelenggarakan pendidikan dan pelatihan pembuatan robot secara reguler, membina ekstrakurikuler robotika, serta mengadakan workshop dan seminar robotika.
    \item \textbf{Bidang Produksi Robotika}: Memproduksi kit edukasi robotika, kit praktikum berbasis komputer, serta kit otomatisasi untuk pesanan umum.
    \item \textbf{Bidang Media dan Publikasi Robotika}: Menulis buku-buku robotika, mempublikasikan hasil riset, membuat modul belajar robotika, dan memberikan edukasi robotika kepada masyarakat melalui platform digital.
    \item \textbf{Bidang Sosial dan Edukasi Masyarakat}: Membangun komunitas-komunitas robotika di kampus dan sekolah melalui program KoTaRo.
\end{enumerate}
